\section{Topics, QoS, Retained Messages und Last Will}

\begin{frame}{Topics, QoS, Retained Messages und Last Will}
  \begin{itemize}
\item \alert{Topics sind hierarchisch}:
	\begin{itemize}
\item \texttt{home/kitchen/temp},
\item \texttt{factory/line1/motor/rpm}
	\end{itemize}
\item \alert{Wildcards} (häufig):
	\begin{itemize}
\item \texttt{+} = eine Ebene (\texttt{home/+/temp})
\item \texttt{\#} = mehrere Ebenen (\texttt{home/\#})
	\end{itemize}
\item \alert{QoS (Quality of Service)} = Zustellgarantie-Stufe:

  \vspace{0.5ex}
  \begin{tabular}{lll}
    \toprule
    \alert{QoS} & \alert{Bezeichnung} & \alert{Bedeutung} \\
    \midrule
    0 & At most once  & Am schnellsten, keine Wiederholungen; Nachricht kann verloren gehen \\
    1 & At least once & Wiederholt bis ACK; Duplikate möglich \\
    2 & Exactly once  & Höchste Zuverlässigkeit; mehr Overhead/Handshakes \\
    \bottomrule
  \end{tabular}

  \vspace{0.5ex}
\item \alert{Retained Messages} (häufig verwendet)
  \begin{itemize}
\item Broker speichert den letzten Wert pro Topic
  \end{itemize}

  \vspace{0.5ex}
\item \alert{Last Will}
  \begin{itemize}
\item Der Broker publiziert eine „Client ist ausgefallen“-Nachricht, wenn ein Client unerwartet die Verbindung verliert
  \end{itemize}
  \end{itemize}
\end{frame}
